\documentclass[10pt,a4paper]{amsart}
\usepackage[margin=19mm]{geometry}
\usepackage{amsmath,amssymb,mathtools}
\usepackage[scr=rsfs,cal=euler]{mathalfa}
\usepackage{xfrac}
\usepackage[unicode,hidelinks,bookmarks=false]{hyperref}
\newcommand{\ie}{i.e.\ }
\newcommand{\cf}{cf.\ }
\newcommand{\resp}{respectively\ }
\newcommand{\Subsection}{\subsection}
\providecommand{\nobreakdash}{\nobreak\hbox{-}}
\setlength{\emergencystretch}{2em}
\multlinegap0pt
\usepackage{bm}
\usepackage{bbm}
\usepackage{dsfont}
\usepackage{xspace}
\usepackage{cancel}
\usepackage{mathtools}
\usepackage{amsthm}
\makeatletter
\@ifundefined{theorem}{\newtheorem{theorem}{Theorem}}{}
\@ifundefined{lemma}{\newtheorem{lemma}[theorem]{Lemma}}{}
\makeatother
\newcommand{\BB}{\mathcal{B}}
\newcommand{\Ball}{\mathsf{B}}
\newcommand{\RRR}{\mathbb{R}}
\newcommand{\cl}{\mathsf{cl}}
\newcommand{\di}{\mathsf{d}}
\title[On the infinite powers of large zero-dimensional metrizable ]{On the infinite powers of large zero-dimensional metrizable spaces}
\author{Andrea Medini}
\date{Source: arXiv:2501.07253v2 (CC BY 4.0)}
\begin{document}
\maketitle

\noindent\emph{Source and licence.} On the infinite powers of large zero-dimensional metrizable spaces. Version arXiv:2501.07253v2 (CC BY 4.0), distributed under the Creative Commons Attribution 4.0 International licence, as recorded in the arXiv licence metadata for this exact version and shown on the abstract page https://arxiv.org/abs/2501.07253v2. Full rights notice: licenses/arxiv-2501.07253-CC-BY-4.0.txt.\par\smallskip

\noindent\textit{Editorial scope.} Counted unit: the Lemma whose source label is \texttt{lemma\_pseudocompact} in the article (source0.tex lines 117--119, source label \texttt{lemma\_pseudocompact}), delivered together with the article's own complete original proof of it (source0.tex lines 120--136, one \texttt{proof} environment of 17 lines). No mathematics is written, repaired or re-derived: statement and proof are the article's own text line for line. Required context retained uncounted: none. Every definition, display and label of those lines is retained unchanged. Omitted from this package and not redistributed: 1--116, 137--213. Nothing inside a retained excerpt is omitted or altered: every retained body is delivered exactly as the article's lines are written, so no omission receipt of any kind accompanies this package. Citations: the retained excerpts carry no citation command at all, so no bibliography entry of the article is transcribed and no reference is redistributed. Reference adaptation: none was needed and none was made; every reference inside the delivered unit resolves inside it, so no unresolved reference or citation is produced. Printed slips kept literal: no printed word, symbol, hypothesis or number of the counted statement or of its proof was altered. Layout and class: the article's own class is \texttt{amsart} with packages including amssymb, amsmath, mathrsfs, verbatim, url, mathabx, xy; the wrapper is the lane's pinned \texttt{amsart} editorial wrapper at 10pt on A4, which restates the theorem-like environments the retained text needs and adds the article's own macro definitions that the retained text needs (\texttt{\textbackslash BB}, \texttt{\textbackslash Ball}, \texttt{\textbackslash RRR}, \texttt{\textbackslash cl}, \texttt{\textbackslash di}), with extra wrapper packages (bm, bbm, dsfont, xspace, cancel, mathtools). The delivered numbering is the unit's own. Assets: no retained span contains a figure environment, a graphics-inclusion command, a PDF-page inclusion, a table or a TikZ picture; no image file is copied into this package, the package declares no asset dependency and no image file is redistributed with it. Licence: version arXiv:2501.07253v2 of this article is distributed under the Creative Commons Attribution 4.0 International licence, as recorded in the arXiv licence metadata for this exact version and shown on the abstract page https://arxiv.org/abs/2501.07253v2. A full rights notice is delivered at licenses/arxiv-2501.07253-CC-BY-4.0.txt. Characters: every retained excerpt is pure ASCII, so the delivered engine, XeLaTeX with its default Latin Modern text font, reports no missing character.
\begin{lemma}\label{lemma_pseudocompact}
Let $X$ be a space. Assume that $X$ has a $\pi$-base $\BB$ consisting of clopen sets, and that there exists an unbounded continuous function $f:X\longrightarrow\RRR$. Then $X$ can be partitioned into infinitely many clopen sets.
\end{lemma}

\begin{proof}
Fix a metric $\di$ on $\RRR$. Throughout this proof, we will use $\cl$ to denote closure in $\RRR$. It is a simple exercise to construct $D=\{d_n:n\in\omega\}\subseteq f[X]$ and open subsets $U_n$ of $\RRR$ for $n\in\omega$ such that the following conditions are satisfied:
\begin{itemize}
\item $D$ is a closed subset of $\RRR$,
\item $d_n\in U_n$ for each $n$,
\item $U_m\cap U_n=\varnothing$ whenever $m\neq n$.
\end{itemize}
Then set $V_n=\Ball(d_n,\varepsilon_n)$ for $n\in\omega$, where the $\varepsilon_n$ are such that $0<\varepsilon_n\leq 2^{-n}$ and $\cl(V_n)\subseteq U_n$.

Next, we will show that $V=\bigcup_{n\in\omega}\cl(V_n)$ is closed in $\RRR$. Pick $x\notin V$. Choose $N\in\omega$ such that $2^{-N}<\di(x,D)$, then set $W=\Ball(x,2^{-(N+1)})$. We claim that $W\cap V_n=\varnothing$ for every $n\geq N+1$. Otherwise, for an element $z$ of such an intersection, we would have
$$
\di(x,d_n)\leq\di(x,z)+\di(z,d_n)\leq 2^{-(N+1)}+2^{-(N+1)}=2^{-N}<\di(x,D),
$$
which is a contradiction. So $W\setminus(\cl(V_0)\cup\cdots\cup\cl(V_N))$ is an open neighborhood of $x$ that is disjoint from $V$.

Finally, fix $B_n\in\BB$ for $n\in\omega$ so that each $B_n\subseteq f^{-1}[V_n]$. To conclude the proof, we will show that $B=\bigcup_{n\in\omega}B_n$ is closed. Pick $x\notin B$. If $x\in f^{-1}[U_n]$ for some $n\in\omega$, then $f^{-1}[U_n]\setminus B_n$ is an open neighborhood of $x$ that is disjoint from $B$. Now assume that $x\notin\bigcup_{n\in\omega} f^{-1}[U_n]$. Then $y=f(x)\notin V$, so we can find an open neighborhood $W$ of $y$ that is disjoint from $V$. It is clear that $f^{-1}[W]$ is an open neighborhood of $x$ that is disjoint from $B$.
\end{proof}

\end{document}
